\documentclass{subfiles}
\begin{document}
\chapter{Alternador trifásico}

\info{Contenido}{Alternadores, frecuencia del alternador,
  sistemas trifásicos, conexiones en estrella y triángulo, potencias
  y corrección del factor de potencia.}

\section{El alternador}
Un alternador es una máquina eléctrica rotativa que transforma
energía mecánica en energía eléctrica de corriente alterna. Su
principio de funcionamiento se basa en la ley de Faraday: cuando un
conductor se mueve dentro de un campo magnético variable, se induce
en él una fuerza electromotriz (fem).

\subsection{Frecuencia del alternador}
La frecuencia de la tensión generada depende de la velocidad de giro
del rotor y del número de pares de polos del inductor. La relación
matemática que los vincula es:

\begin{equation}
  f = \frac{n \cdot p}{60}
\end{equation}

Donde:
\begin{itemize}
  \item $f$ = frecuencia de la tensión generada, en Hertz (\unit{Hz}).
  \item $n$ = velocidad de rotación, en revoluciones por minuto (\unit{rpm}).
  \item $p$ = número de pares de polos.
\end{itemize}

\section{Sistemas trifásicos}
Un sistema trifásico es un conjunto de tres tensiones alternas de
igual amplitud y frecuencia, desfasadas entre sí $120^\circ$. Es el
sistema estándar para la generación y distribución de energía
eléctrica debido a su eficiencia y constancia en la potencia
instantánea entregada.

En un sistema trifásico se distinguen dos tipos de tensiones y dos
tipos de corrientes:
\begin{itemize}
  \item \textbf{Tensión simple} ($V_S$): Es la tensión entre una fase
    y el neutro.
  \item \textbf{Tensión compuesta} ($V_C$): Es la tensión entre dos
    fases cualesquiera.
  \item \textbf{Corriente simple} ($I_S$): Es la corriente que circula
    por la carga en cada fase.
  \item \textbf{Corriente compuesta} ($I_C$): Es la corriente que
    circula por las líneas de alimentación.
\end{itemize}

\section{Conexiones de un sistema trifásico}
Las cargas trifásicas se pueden conectar principalmente de dos
formas: en estrella (Y) o en triángulo ($\Delta$).

\subsection{Conexión estrella (Y)}
En la conexión estrella, un extremo de cada impedancia se une a un
punto común llamado neutro. En este sistema se cumple:
\begin{equation}
  V_C = \sqrt{3} \cdot V_S \approx 1{,}732 \cdot V_S
\end{equation}

\subsection{\texorpdfstring{Conexión triángulo ($\Delta$)}{Conexión triángulo (Delta)}}
En la conexión triángulo, las impedancias se cierran formando un
lazo, y las líneas se conectan a los vértices. En este sistema se
cumple:
\begin{equation}
  I_C = \sqrt{3} \cdot I_S \approx 1{,}732 \cdot I_S
\end{equation}

\section{Potencias en sistemas trifásicos}
En un sistema trifásico equilibrado, las potencias totales se
calculan utilizando los valores compuestos (de línea):

\begin{equation}
  P = \sqrt{3} \cdot V_C \cdot I_C \cdot \cos(\varphi)
\end{equation}
\begin{equation}
  Q = \sqrt{3} \cdot V_C \cdot I_C \cdot \sin(\varphi)
\end{equation}
\begin{equation}
  S = \sqrt{3} \cdot V_C \cdot I_C
\end{equation}

Además, la potencia aparente se relaciona con las potencias activa
y reactiva mediante:
\begin{equation}
  S = \sqrt{P^2 + Q^2}
\end{equation}

\subsection{Corrección del factor de potencia}
Para corregir el factor de potencia conectando un banco de
condensadores en paralelo a la carga, la potencia reactiva del banco
se calcula como:
\begin{equation}
  Q_C = P \left( \tan(\varphi_1) - \tan(\varphi_2) \right)
\end{equation}


\section{Problemas resueltos}
\actividad{Resolver los siguientes problemas}
\begin{enumerate}
  \item Determinar la frecuencia que produce un alternador que gira a
    una velocidad de $n=\qty{1500}{rpm}$ si éste posee dos pares de
    polos.
  \item Determinar la frecuencia que produce un alternador que gira a
    una velocidad de $n=\qty{3500}{rpm}$ si posee 4 pares de polos.
  \item Determinar el número de pares de polos que posee un generador
    que gira a una velocidad de $n=\qty{4500}{rpm}$ y genera una
    frecuencia de \qty{300}{Hz}.
  \item Determinar el número de pares de polos que posee un generador
    que gira a una velocidad de $n=\qty{3000}{rpm}$ y genera una
    frecuencia de \qty{100}{Hz}.
  \item ¿A qué velocidad deberá girar un alternador con cuatro pares
    de polos para producir una frecuencia de $f=\qty{50}{Hz}$?
  \item Determinar la tensión compuesta que corresponde a un sistema
    trifásico que posee una tensión simple de $V_S=\qty{223}{V}$.
  \item Determinar la tensión simple que corresponde a un sistema
    trifásico que posee una tensión compuesta de $V_C=\qty{440}{V}$.
  \item Un motor trifásico posee sus bobinas conectadas en estrella.
    Determinar la corriente eléctrica que absorberá de la línea si al
    conectarlo a una red con una tensión entre fases de
    $V_C=\qty{400}{V}$ desarrolla una potencia de $P=\qty{10}{kW}$ con
    un factor de potencia de $\cos(\varphi)=0{,}8$. Averiguar la
    potencia reactiva $Q$ y aparente $S$ del mismo.
  \item Se conectan en estrella tres bobinas iguales a una red
    trifásica de $V_C=\qty{220}{V}$, $f=\qty{50}{Hz}$. Cada una de las
    mismas posee $R=\qty{12}{\Omega}$ de resistencia, y
    $X_L=\qty{20}{\Omega}$ de reactancia inductiva. Calcular la
    impedancia $Z$, corrientes simples, factor de potencia
    $\cos(\varphi)$, potencia activa $P$, potencia reactiva $Q$ y
    potencia aparente $S$.
  \item Se desean conectar a una red trifásica con neutro y una
    tensión entre fases de $V_C=\qty{400}{V}$, 30 lámparas
    fluorescentes de $P=\qty{20}{W}$, $\qty{230}{V}$,
    $\cos(\varphi)=0{,}6$.
    \begin{itemize}
      \item ¿Cómo deben conectarse para que la carga esté equilibrada?
      \item Averiguar las corrientes simples y compuestas.
      \item Calcular la potencia del conjunto y por fase.
    \end{itemize}
  \item Un motor trifásico posee sus bobinas conectadas en triángulo.
    Determinar la corriente eléctrica que absorberá de la línea si al
    conectarlo a una red con tensión entre fases de
    $V_C=\qty{380}{V}$, se desarrolla una potencia de
    $P=\qty{15}{kW}$, con $\cos(\varphi)=0{,}7$. Averiguar la potencia
    reactiva $Q$ y aparente $S$.
  \item Se conectan en triángulo bobinas de $R=\qty{10}{\Omega}$ y
    $X_L=\qty{10}{\Omega}$ a una red trifásica de $V_C=\qty{400}{V}$,
    $f=\qty{50}{Hz}$. Calcular $I_C$, $I_S$, $\cos(\varphi)$, $P$,
    $Q$, $S$.
  \item Se desean conectar 54 lámparas incandescentes de
    $\qty{100}{W}$, $\qty{220}{V}$, a una red trifásica con una
    tensión entre fases de $V_C=\qty{220}{V}$. ¿Cómo deben conectarse
    y qué consumo poseen?
  \item En un sistema trifásico con carga equilibrada se mide una
    intensidad de línea de $I_C=\qty{30}{A}$, con un factor de
    potencia de $\cos(\varphi)=0{,}75$. Si la tensión entre fases es
    de $V_C=\qty{380}{V}$. Averiguar las potencias $Q$, $P$, $S$.
  \item En un sistema trifásico con carga equilibrada de tres hilos
    se mide una potencia de $P=\qty{36}{kW}$ y una intensidad de
    $I_C=\qty{97{,}4}{A}$. Si la tensión entre fases es de
    $V_C=\qty{380}{V}$, averiguar el factor de potencia
    $\cos(\varphi)$ de la carga.
  \item Un aparato de calefacción trifásico consta de tres
    resistencias con $R=\qty{10}{\Omega}$ conectadas en estrella.
    Determinar la potencia que consumirán cuando se aplique
    $V_C=\qty{230}{V}$ entre fases, así como la corriente simple $I_S$
    y compuesta $I_C$.
  \item Si conectamos en triángulo las resistencias del problema
    anterior, ¿qué valores tomarán $I_C$ y $I_S$?
  \item Un motor trifásico de $P=\qty{3990}{W}$ y
    $\cos(\varphi)=0{,}65$ se conecta a una red de
    $V_C=\qty{380}{V}$, $f=\qty{50}{Hz}$. Se trata de averiguar la
    corriente de línea y de cada fase del motor cuando está conectado
    en triángulo, así como su potencia activa $P$, reactiva $Q$ y
    aparente $S$. Si cada una de las fases del motor se puede
    considerar como una impedancia $Z$ que posee una reactancia
    inductiva $X_L$ y una resistencia $R$, determinar el valor de las
    mismas.
  \item Se conectan en triángulo tres bobinas iguales de
    $R=\qty{16}{\Omega}$ y $X_L=\qty{10}{\Omega}$ de reactancia
    inductiva. Si se conectan en un sistema trifásico de
    $V_C=\qty{240}{V}$ entre fases y $f=\qty{50}{Hz}$. Determinar:
    \begin{itemize}
      \item Las corrientes simples $I_S$ y compuestas $I_C$.
      \item El factor de potencia $\cos(\varphi)$.
      \item La potencia activa $P$ y reactiva $Q$.
    \end{itemize}
  \item Una red trifásica alimenta tres motores monofásicos de
    inducción de $P=\qty{5}{CV}$ ($\qty{1}{CV}=\qty{736}{W}$),
    $\cos(\varphi)=0{,}78$, $V_C=\qty{220}{V}$, cada uno conectado
    entre fase y neutro. Determinar:
    \begin{itemize}
      \item Las corrientes simples $I_S$ y compuestas $I_C$.
      \item La potencia reactiva que deberá aportar un banco de
        condensadores para corregir el factor de potencia a
        $\cos(\varphi)=0{,}9$.
    \end{itemize}
  \item Una empresa demanda una potencia aparente $S=\qty{700}{kVA}$
    a $V_C=\qty{10}{kV}$ en corriente alterna trifásica. Las lecturas
    de consumo en dos meses son: contador de energía activa
    $P_h=\qty{205000}{kWh}$ y de energía reactiva
    $Q_h=\qty{150000}{kvar\cdot h}$. Calcular:
    \begin{itemize}
      \item El $\cos(\varphi)$ medio en dicho período de facturación.
      \item Las intensidades $I_C$ e $I_S$.
      \item La potencia reactiva de un banco de condensadores
        necesaria para corregir el factor de potencia a
        $\cos(\varphi)=0{,}93$.
      \item El porcentaje de reducción de la corriente simple $I_S$
        una vez conectado el banco de condensadores.
    \end{itemize}
  \item Una instalación industrial de $\qty{50}{kW}$, con un factor
    de potencia de $\cos(\varphi)=0{,}65$, se alimenta a través de un
    transformador trifásico, con una tensión en el primario entre
    fases de $V_{C1}=\qty{24}{kV}$ y de $V_{C2}=\qty{400}{V}$ en el
    secundario. Averiguar:
    \begin{itemize}
      \item La potencia nominal del transformador en kVA, así como la
        corriente por el primario y el secundario.
      \item Las nuevas características del transformador si se
        corrige el factor de potencia a $\cos(\varphi)=0{,}98$
        mediante una batería automática de condensadores conectada en
        el lado de baja tensión.
    \end{itemize}
\end{enumerate}
\end{document}